% Explicit modern presentation. Dummett's original denumerable matrix remains distinct.
\calculussheet{sheet:hilbert-lc}{Hilbert Presentation of LC}{Language: atoms and $\bot,\land,\lor,\to$, with $\neg A\defeq A\to\bot$; judgment: $\Gamma\derives_{\mathsf H_{\mathsf{LC}}}A$ by a finite sequence of premise lines, instances of the displayed schemas, and modus ponens. $A,B,C$ range over all formulas of this propositional language. Uniform substitution instantiates schemas, not arbitrary premise lines. Equality and quantifiers: absent.}
\noindent\emph{Source distinction:} this fixed full-language Hilbert base plus prelinearity is a modern presentation of the theorem set profiled in \cite{Dummett1959}; the finite-chain bridge below is not a reconstruction of that paper's particular denumerable matrix. The local frame proof is Theorem~\ref{thm:p08-mi-intermediate-adequacy}.
\hbrule{rule:hilbert-lc-a1}{A1}{}{A\to(B\to A)}{\none}
\hbrule{rule:hilbert-lc-a2}{A2}{}{(A\to(B\to C))\to((A\to B)\to(A\to C))}{\none}
\hbrule{rule:hilbert-lc-and-i}{Conjunction}{}{A\to(B\to A\land B)}{\none}
\hbrule{rule:hilbert-lc-and-e1}{Projection 1}{}{A\land B\to A}{\none}
\hbrule{rule:hilbert-lc-and-e2}{Projection 2}{}{A\land B\to B}{\none}
\hbrule{rule:hilbert-lc-or-i1}{Injection 1}{}{A\to A\lor B}{\none}
\hbrule{rule:hilbert-lc-or-i2}{Injection 2}{}{B\to A\lor B}{\none}
\hbrule{rule:hilbert-lc-or-e}{Cases}{}{(A\to C)\to((B\to C)\to((A\lor B)\to C))}{\none}
\hbrule{rule:hilbert-lc-bot}{Falsity}{}{\bot\to A}{\none}
\hbrule{rule:hilbert-lc-prelinearity}{Prelinearity}{}{(A\to B)\lor(B\to A)}{\none}
\hbrule{rule:hilbert-lc-mp}{MP}{A\to B\qquad A}{B}{both premise lines precede this conclusion line in the sequence}
